% Copyright (c) 2026 Benjamin Stanley Frohman.
% Licensed under the Apache License, Version 2.0.
% REVIEW DRAFT --- not a DOI assignment, upload, or publication.
\documentclass[11pt,a4paper]{article}

\usepackage[margin=1.05in]{geometry}
\usepackage{amsmath,amssymb,amsthm,mathtools}
\usepackage{microtype}
\usepackage{enumitem}
\usepackage{booktabs}
\usepackage{longtable}
\usepackage{hyperref}
\usepackage[T1]{fontenc}
\usepackage{lmodern}
\usepackage{setspace}
\usepackage{xcolor}
\usepackage{xurl}
\Urlmuskip=0mu plus 3mu
\setstretch{1.08}

\hypersetup{colorlinks=true,linkcolor=black,citecolor=black,urlcolor=blue}

\newcommand{\Q}{\mathbb{Q}}
\newcommand{\C}{\mathbb{C}}
\newcommand{\lean}[1]{\texttt{\detokenize{#1}}}

\newtheorem{theorem}{Theorem}
\newtheorem{lemma}[theorem]{Lemma}
\newtheorem{corollary}[theorem]{Corollary}
\theoremstyle{definition}
\newtheorem{definition}[theorem]{Definition}
\newtheorem{remark}[theorem]{Remark}

\title{\textbf{Binary Jacobian Quotient Representatives\\
and H/G Common-Zero Exclusion}\\[0.35em]
\large Lean Stage~3 Algebraic Certificates\\[0.4em]
\normalsize\textcolor{red}{\textbf{REVIEW DRAFT}} --- not peer-reviewed; no DOI assigned; not uploaded}
\author{Frohman, Benjamin Stanley\\
\small ORCID: \href{https://orcid.org/0009-0006-7068-3718}{0009-0006-7068-3718}\\
\small Email: \texttt{frohmanbenjamin@gmail.com}\\
\small Independent Researcher}
\date{6~November~2026 (America/Chicago)\\
\small Source pin: \texttt{e56be3d0e848c2098f1d9057dccec324f34bae7f}\\
\small Repository: \url{https://github.com/BenFrohman/HODGE}\\
\small Pull request: \url{https://github.com/BenFrohman/HODGE/pull/38}}

\begin{document}
\maketitle

\begin{center}
\fbox{\parbox{0.92\textwidth}{\centering\small
\textbf{Status.} This manuscript is a \textbf{REVIEW DRAFT} companion to the
Stage~1--2 algebraic notes. It invents no DOI, performs no upload or
publication, and does not claim peer review, novelty, or a Clay/Hodge close.
License of the cited Lean sources remains Apache-2.0; this note does not
relicense them. Author metadata uses the exact form
\texttt{Frohman, Benjamin Stanley}.}}
\end{center}

\begin{abstract}
We record thirteen Lean~4 theorems in
\lean{Hodge.ExplicitJacobianQuotient} at commit
\path{e56be3d0e848c2098f1d9057dccec324f34bae7f}. Over the actual binary
polynomial ring $K[x,y]$, with $f_k=kx^5y+y^6$, the file identifies the
Jacobian ideal generated by the two formal partial derivatives, proves the
corresponding quotient relations (including $x^4y=0$ and $y^6=0$ when $k\neq0$),
identifies the section Wronskian determinant $W$ with the Stage~2 coefficient
formulas, proves the four-monomial identity
$[HW]=u_0 y^5+u_1 xy^4+u_2 x^2y^3+u_3 x^3y^2$ in the quotient at every modulus
root, proves $G^2-H^3=f_k$ in the polynomial ring from the six existing
coefficient relations, and proves that any common evaluation zero of the actual
$H$ and $G$ forces the origin. Combined with Stages~1--2 this yields a local
audit of $65$ incremental theorem declarations under only
\lean{propext}, \lean{Classical.choice}, and \lean{Quot.sound}.
This is not a Sylvester-resultant identification, not a quotient basis or
dimension theorem, and not a geometric Hodge/residue/cycle-class proof.
\end{abstract}

\tableofcontents

\section{Scope and relation to Stages~1--2}

This note is a third coherent algebraic paper for a distinct substantive
result: the actual binary Jacobian-quotient representatives and the
projective common-zero exclusion for the section forms $H,G$. Supporting
lemmas that are solely scaffolding for that result are bundled here; Stages~1
and~2 remain separate papers for the determinant/derivative identities and the
B\'ezout nonvanishing certificates~\cite{HODGERepo}.

\paragraph{Source.}
Module \path{Hodge/ExplicitJacobianQuotient.lean} at
\path{e56be3d0e848c2098f1d9057dccec324f34bae7f}
(PR~\url{https://github.com/BenFrohman/HODGE/pull/38}).
Byte SHA-256 of the Lean source:
\path{0ff4a38ce58504ad4e4b53f327665e690895c795db34ebcbbd99159ce97e6ae5}
(matches \path{stage3-manifest.json}).

\paragraph{What this note proves.}
Exactly the thirteen theorems listed in Section~\ref{sec:mapping}, with
hypotheses taken from the Lean declarations (no hidden assumptions).

\paragraph{What this note does not prove.}
\begin{itemize}[leftmargin=1.4em,itemsep=0.25em]
\item Identification of mathlib's Sylvester resultant, or a Lean check of the
      tenth Python numeric resultant/B\'ezout certificate.
\item Canonical normal forms, basis independence, quotient dimension, or that
      a nonzero representative coefficient implies a nonzero quotient class.
\item Reducedness/saturation of geometric ideals, schemes, Betti/Hodge
      comparison, residue-cycle theorems, Galois/Fourier descent, or any
      general/integral Hodge conjecture statement.
\end{itemize}
Terminology: the binary determinant $W$ below is the Wronskian of the section
forms $H,G$ (formal partials). It is \emph{not} Jacobi's identity for Lie
brackets, and it is not a claim that $W$ is a cohomology class.

\section{Standing definitions (exact Lean objects)}

Throughout, $K$ is a commutative ring, a field, or a characteristic-zero field
with a $\Q$-algebra structure, as required by each theorem.

\begin{definition}[Binary ring and $f_k$]
\lean{Binary K}${}={}$\lean{MvPolynomial (Fin 2) K}, with $x=X(0)$, $y=X(1)$, and
\begin{equation}
f_k \;=\; k\,x^5 y + y^6
\qquad(\lean{binarySextic}\ k).
\end{equation}
This is the scaled section polynomial; the unscaled $x^5y+y^6$ is not substituted.
\end{definition}

\begin{definition}[Jacobian ideal and quotient]
\begin{align}
J(k) &\;=\; \bigl(\partial_x f_k,\;\partial_y f_k\bigr)
\qquad(\lean{jacobianIdeal}\ k),\\
\operatorname{Jac}(k) &\;=\; K[x,y]/J(k)
\qquad(\lean{Jacobian}\ k).
\end{align}
Write $\overline{x},\overline{y}$ for the images of $x,y$
(\lean{qx},\lean{qy}), and $c\mapsto\overline{c}$ for the constant map
(\lean{constants}).
\end{definition}

\begin{definition}[Section forms]
With Stage~2/ExplicitCurve coefficients $b=\mathrm{sb}(k)$, \ldots,
$g=\mathrm{sg}(k)$,
\begin{align}
H &\;=\; x^2 + b\,xy + c\,y^2,\\
G &\;=\; x^3 + d\,x^2y + e\,xy^2 + g\,y^3,\\
W &\;=\; \partial_x H\cdot\partial_y G - \partial_y H\cdot\partial_x G.
\end{align}
These are \lean{sectionH}, \lean{sectionG}, \lean{sectionW}.
The Stage~2 scalars $w_i$, $u_i$ are exactly the certificates
\path{w0},\ldots,\path{u3} from \path{ExplicitResidueCertificates}.
\end{definition}

\begin{definition}[Modulus-root hypothesis]
\begin{equation}\label{eq:mod}
k^{24}+66081312000\,k^{18}-4811512860000\,k^{12}
+1167566400000\,k^{6}+583200000=0.
\end{equation}
Irreducibility of this polynomial over $\Q$ is neither assumed nor proved.
\end{definition}

\section{Main results}

\begin{theorem}[Partial derivatives of $f_k$]\label{thm:partials}
For a commutative ring $K$ and $k\in K$,
\begin{equation}
\partial_x f_k = 5k\,x^4 y,\qquad
\partial_y f_k = k\,x^5 + 6y^5.
\end{equation}
\end{theorem}
\begin{proof}[Lean]
\lean{partial_x}, \lean{partial_y}: unfold and ring/simp.
\end{proof}

\begin{theorem}[Generators vanish in the quotient]\label{thm:partial-classes}
In $\operatorname{Jac}(k)$,
\begin{equation}
\overline{5k}\,\overline{x}^4\overline{y}=0,\qquad
\overline{k}\,\overline{x}^5 + 6\,\overline{y}^5=0.
\end{equation}
In particular the second relation is $\overline{k}\,\overline{x}^5=-6\overline{y}^5$,
not $\overline{x}^5=-6\overline{y}^5$.
\end{theorem}
\begin{proof}[Lean]
\lean{partial_x_class}, \lean{partial_y_class}: membership of the generators
in $J(k)$ and the quotient zero criterion.
\end{proof}

\begin{lemma}[Constant inverse]\label{lem:cinv}
If $K$ is a field and $a\in K\setminus\{0\}$, then
$\overline{a^{-1}}\,\overline{a}=1$ in $\operatorname{Jac}(k)$.
\end{lemma}
\begin{proof}[Lean]
\lean{constant_inverse}.
\end{proof}

\begin{theorem}[Monomial vanishings for $k\neq0$]\label{thm:vanishing}
Assume $K$ is a field of characteristic zero and $k\neq0$. Then in
$\operatorname{Jac}(k)$,
\begin{equation}
\overline{x}^4\overline{y}=0\qquad\text{and}\qquad
\overline{y}^6=0.
\end{equation}
\end{theorem}
\begin{proof}[Lean]
\lean{x4y_class}: invert $\overline{5k}$ using Lemma~\ref{lem:cinv} and
Theorem~\ref{thm:partial-classes}.
\lean{y6_class}: combine $\overline{y}\cdot(\partial_y\text{-class})$ with
$\overline{k}\,\overline{x}\cdot(\overline{x}^4\overline{y}=0)$ and invert $6$.
\end{proof}

\begin{theorem}[Wronskian equals Stage~2 $W$-coefficients]\label{thm:Wpoly}
Assume $K$ is a characteristic-zero field with a $\Q$-algebra structure. Then
\begin{equation}
W = C(w_3)\,x^3+C(w_2)\,x^2y+C(w_1)\,xy^2+C(w_0)\,y^3
\end{equation}
in $K[x,y]$, where each $w_i\in K$ is the Stage~2 scalar coefficient and
$C\colon K\to K[x,y]$ embeds constants. The quotient image in
$\operatorname{Jac}(k)$ retains bars only on the quotient representatives:
\begin{equation}
\overline{W}=
\overline{C(w_3)}\,\overline{x}^3+
\overline{C(w_2)}\,\overline{x}^2\overline{y}+
\overline{C(w_1)}\,\overline{x}\,\overline{y}^2+
\overline{C(w_0)}\,\overline{y}^3.
\end{equation}
\end{theorem}
\begin{proof}[Lean]
\lean{wronskian_polynomial} applies Stage~1
\lean{binary_wronskian} at the actual section coefficients;
\lean{wronskian_class} pushes forward to the quotient.
\end{proof}

\begin{theorem}[$H$ class]\label{thm:Hclass}
Under the same field/$\Q$-algebra hypotheses,
\begin{equation}
\overline{H}=\overline{x}^2+\overline{b}\,\overline{x}\,\overline{y}+\overline{c}\,\overline{y}^2.
\end{equation}
\end{theorem}
\begin{proof}[Lean]
\lean{sectionH_class}.
\end{proof}

\begin{theorem}[{Four-monomial $[HW]$ identity}]\label{thm:HW}
Assume $K$ is a characteristic-zero field with a $\Q$-algebra structure and
that $k$ satisfies~\eqref{eq:mod}. Then in $\operatorname{Jac}(k)$,
\begin{equation}
\overline{HW}
=
\overline{u_0}\,\overline{y}^5
+\overline{u_1}\,\overline{x}\,\overline{y}^4
+\overline{u_2}\,\overline{x}^2\overline{y}^3
+\overline{u_3}\,\overline{x}^3\overline{y}^2,
\end{equation}
with the Stage~2 definitions
$u_0=cw_0-6$, $u_1=bw_0+cw_1$, $u_2=w_0+bw_1+cw_2$, $u_3=w_1+bw_2+ck$.
The constant $-6$ arises from the $k x^5$ term together with the actual
$\partial_y$ relation; it is not obtained by setting $k=1$.
\end{theorem}
\begin{proof}[Lean]
\lean{section_HW_class}: use $w_3=k$ via \lean{ExplicitCurve.rel_h5}, multiply
the $H$ and $W$ classes, and reduce with Theorems~\ref{thm:vanishing}
and~\ref{thm:partial-classes} ($k\neq0$ from
\lean{section_root_ne_zero}).
\end{proof}

\begin{theorem}[Section identity in the polynomial ring]\label{thm:identity}
Under the same hypotheses as Theorem~\ref{thm:HW},
\begin{equation}
G^2-H^3 = f_k
\end{equation}
as an equality in $K[x,y]$, derived from the six existing checked coefficient
relations (\path{rel_h5} through \path{rel_h0}) and~\eqref{eq:mod}, not
postulated.
\end{theorem}
\begin{proof}[Lean]
\lean{section_identity_polynomial}: push the six scalar relations into
constants and apply \lean{ExplicitCurve.section_identity}.
\end{proof}

\begin{theorem}[Common-zero exclusion]\label{thm:common}
Under the same hypotheses as Theorem~\ref{thm:HW}, if $a,b\in K$ satisfy
\begin{equation}
H(a,b)=0\qquad\text{and}\qquad G(a,b)=0,
\end{equation}
then $a=0$ and $b=0$.
\end{theorem}
\begin{proof}[Lean outline]
\lean{section_common_zero}. Differentiate Theorem~\ref{thm:identity}: at a
common zero of $H$ and $G$, both partials of $f_k$ vanish, so
\begin{equation}
5k a^4 b=0,\qquad k a^5+6b^5=0.
\end{equation}
From~\eqref{eq:mod} one has $k\neq0$, hence $a=0$ or $b=0$; the second equation
forces the remaining coordinate to vanish. No resultant API, coprimality
premise, or external geometric theorem is assumed. This discharges the
projective common-zero input supplied classically by the tenth Python
certificate, without identifying Sylvester resultants.
\end{proof}

\section{Theorem--declaration mapping}
\label{sec:mapping}

\begin{center}
\begin{longtable}{@{}llp{6.0cm}@{}}
\toprule
\# & Lean declaration & Content \\
\midrule
1 & \lean{partial_x} & Thm.~\ref{thm:partials}: $\partial_x f_k$ \\
2 & \lean{partial_y} & Thm.~\ref{thm:partials}: $\partial_y f_k$ \\
3 & \lean{partial_x_class} & Thm.~\ref{thm:partial-classes}: $\overline{5k}\,\overline{x}^4\overline{y}=0$ \\
4 & \lean{partial_y_class} & Thm.~\ref{thm:partial-classes}: $\overline{k}\,\overline{x}^5+6\overline{y}^5=0$ \\
5 & \lean{constant_inverse} & Lem.~\ref{lem:cinv} \\
6 & \lean{x4y_class} & Thm.~\ref{thm:vanishing}: $\overline{x}^4\overline{y}=0$ \\
7 & \lean{y6_class} & Thm.~\ref{thm:vanishing}: $\overline{y}^6=0$ \\
8 & \lean{wronskian_polynomial} & Thm.~\ref{thm:Wpoly}: $W$ coefficients \\
9 & \lean{wronskian_class} & Thm.~\ref{thm:Wpoly}: $\overline{W}$ \\
10 & \lean{sectionH_class} & Thm.~\ref{thm:Hclass} \\
11 & \lean{section_HW_class} & Thm.~\ref{thm:HW}: four-monomial $[HW]$ \\
12 & \lean{section_identity_polynomial} & Thm.~\ref{thm:identity}: $G^2-H^3=f_k$ \\
13 & \lean{section_common_zero} & Thm.~\ref{thm:common}: only origin \\
\bottomrule
\end{longtable}
\end{center}

\noindent
Namespace: \lean{Hodge.ExplicitJacobianQuotient}. Supporting lemmas for this
result are kept in this paper; Stages~1--2 are not re-split here.

\section{Verification evidence (local cached)}

Primary local evidence (this packaging does not re-run Lake):
\begin{itemize}[leftmargin=1.4em]
\item \path{/workspace/hodge-stage3-evidence-20261106/build-and-audit-final.log}:
      \texttt{OK [7086/7086] Built Hodge.ExplicitJacobianQuotient};
      \texttt{PASS: 65 incremental theorem declarations; standard axioms only}.
\item \path{axioms-65.json}: $65$ declarations ($7{+}45{+}13$), each using only
      the standard axioms \lean{propext}, \lean{Classical.choice}, and
      \lean{Quot.sound}.
\item \path{no-sorry-final.log}:
      ``OK: no sorry/admit/native\_decide, no axiom declarations''.
\item \path{full-cached-build.log}: full repository cached build completed.
\item \path{stage3-manifest.json}: commit/source/evidence hashes; scope string
      records ``no basis-independence, nonzero quotient-class or geometric Hodge
      theorem claim''.
\end{itemize}

The independent review typecheck log
\path{/workspace/hodge-stage3-independent-review-typecheck.log} is empty
(zero bytes). At most its process exit status can support that a typecheck
invocation completed; it supplies \emph{no} file-text axiom evidence and is
not cited as an audit. Dedicated axiom/audit evidence is only the items above
(\path{build-and-audit-final.log}, \path{axioms-65.json},
\path{no-sorry-final.log}). Intermediate scratch logs such as
\path{typecheck-2.log} record pre-final errors and are not the passing audit.

Post-merge CI on \texttt{main} reports SUCCESS for all six Lean / Lean
certificate / Paper runs of merges~\#37 and~\#38
(certificate audits: Stage~1--2 PASS~$52$ on merge~\#37;
Stage~1--3 PASS~$65$ with three residue modules on merge~\#38, completed
about 18:48 America/Chicago). Paper CI builds only
\path{paper/hodge_sextic/}; it does not rebuild these formalization notes.
Mathematical source pin remains \path{e56be3d}; merge tip
\path{9289b3c6ae3f2f3ef616493737c113a09c255330}. Manuscript status remains
REVIEW DRAFT regardless of CI.

\section{Toolchain pins and reproducibility}

\begin{center}
\begin{tabular}{@{}ll@{}}
\toprule
Pin & Value \\
\midrule
Lean & \lean{leanprover/lean4:v4.22.0} \\
mathlib input & \lean{v4.22.0} \\
mathlib checkout & \path{79e94a093aff4a60fb1b1f92d9681e407124c2ca} \\
Git commit & \path{e56be3d0e848c2098f1d9057dccec324f34bae7f} \\
PR & \url{https://github.com/BenFrohman/HODGE/pull/38} \\
License & Apache-2.0 (preserved; no relicense) \\
\bottomrule
\end{tabular}
\end{center}

\begin{verbatim}
lake build Hodge.ExplicitJacobianQuotient
python3 scripts/check_residue_axioms.py
python3 scripts/check_no_sorry.py
\end{verbatim}

\section{Limitations}

Theorem~\ref{thm:HW} exhibits a representative of $\overline{HW}$ with
coefficients $u_i$. Nonvanishing of those scalars (Stage~2) does \emph{not}
by itself prove that the corresponding homogeneous components are nonzero as
quotient classes; a basis/independence theorem for the Jacobian quotient is
still required. Theorem~\ref{thm:common} excludes common nonzero zeros of
$H,G$ but does not formalize Proj, schemes, or residue maps.
Villaflor Loyola's residue-cycle theorem~\cite{Villaflor2021} and Griffiths
comparison remain external classical input.

\section*{Author, license, and contact}
Author metadata: \textbf{Frohman, Benjamin Stanley}.
Email: \texttt{frohmanbenjamin@gmail.com}.
ORCID: \href{https://orcid.org/0009-0006-7068-3718}{0009-0006-7068-3718}
(verified from repository \path{CITATION.cff}/\path{NOTICE}/\path{.zenodo.json}).
Copyright (c) 2026 Benjamin Stanley Frohman. Apache License 2.0.
This REVIEW DRAFT assigns no DOI. No Zenodo write was performed for this note.

\bibliographystyle{plain}
\bibliography{references}

\end{document}
